\section{Motivation and Context}
\label{sec:motivation}

Institutional investors play a central role in modern equity markets, holding large positions across a broad universe of securities. Their portfolio decisions therefore have implications not only for individual asset prices but also for the way shocks propagate across securities. In particular, the risk associated with a stock may depend not only on its fundamentals or aggregate level of institutional ownership but also on the composition and structure of its investor base.

This ownership structure becomes especially relevant when positions are crowded. A position can be considered crowded when many institutional investors are exposed to the same security or to similar groups of securities. Crowding is not necessarily problematic under normal market conditions. However, when investors holding overlapping portfolios face common liquidity shocks, binding risk constraints, or deteriorating market conditions, several institutions may attempt to sell the same securities simultaneously. Such trading can generate non-fundamental selling pressure and amplify price movements, particularly when market liquidity is limited. Evidence from mutual-fund fire sales shows that flow-induced institutional trading can exert price pressure on commonly held securities \parencite{coval2007asset}. This mechanism suggests that the identity and characteristics of a stock's owners may contain information not captured by conventional firm-level risk measures. \Textcite{greenwood2011stock} formalize this intuition through the concept of \emph{stock price fragility}, according to which an asset becomes more vulnerable to non-fundamental demand shocks when its ownership is concentrated or its investors face correlated liquidity shocks. Using hedge-fund holdings, \textcite{brown2022crowded} further show that crowded positions are associated with downside tail risk and relatively larger drawdowns during periods of industry distress.

Traditional crowding measures generally summarize institutional ownership using stock-level quantities, such as the percentage of shares held by institutions, the number of institutional holders, or ownership concentration. Although informative, these measures compress a complex ownership structure into a small number of scalar variables. They do not fully represent which institutions hold a security, how similar those institutions' portfolios are, or how the same investors connect different securities through common holdings. Two stocks with comparable levels of institutional ownership may therefore have substantially different underlying ownership structures and, potentially, different degrees of vulnerability to coordinated trading. A network representation provides a natural way to preserve this additional information. Institutional investors and securities can be represented as interconnected nodes, with ownership positions forming the links between them. This representation makes it possible to characterize portfolio overlap, investor similarity, community concentration, centrality, and other dimensions of interconnectedness that conventional crowding measures cannot capture directly. Graph-based machine-learning models can also learn complex and nonlinear patterns from this ownership structure rather than relying exclusively on manually constructed indicators.

Against this background, this thesis investigates whether the topology of institutional ownership contains incremental information about future stock-level downside risk beyond conventional ownership measures and market characteristics. It also examines whether graph-based models can exploit this relational structure more effectively than standard tabular models. The objective is not simply to introduce more complex modeling techniques, but to determine whether their additional complexity produces meaningful improvements in out-of-sample prediction. The resulting predictions are evaluated in terms of both statistical performance and their potential relevance for portfolio construction and investment decisions.